% ============================================================
%  cover.tex — 数学物理方法笔记封面（九层设计）
%  主色调：WaveMain 青蓝 + Gold 暖金
%  用法：\coverimage{figure/封面.jpg} → \makecover
% ============================================================

\makeatletter
\providecommand{\@coverimage}{}
\newcommand{\coverimage}[1]{\renewcommand{\@coverimage}{#1}}
\newcommand{\makecover}{%
\begin{titlepage}
\thispagestyle{empty}

% 第一层：左侧青蓝竖条 + 渐变纹理 + 装饰圆点
\begin{tikzpicture}[remember picture,overlay]
  \fill[WaveMain!90]
    (current page.south west) rectangle
    ([xshift=2.8cm]current page.north west);
  \shade[top color=WaveDark!90, bottom color=WaveMain!90]
    ([xshift=0cm,yshift=-5cm]current page.north west) rectangle
    ([xshift=2.8cm]current page.north west);
  \shade[bottom color=WaveDark!90, top color=WaveMain!90]
    ([xshift=0cm,yshift=5cm]current page.south west) rectangle
    ([xshift=2.8cm]current page.south west);
  \draw[WaveLight!40,line width=1pt]
    ([xshift=2.8cm]current page.south west) --
    ([xshift=2.8cm]current page.north west);
  \draw[WaveLight!20,line width=0.4pt]
    ([xshift=0.4cm]current page.south west) --
    ([xshift=0.4cm]current page.north west);

  \foreach \y/\r/\o in {
    0.08/1.1/0.06, 0.22/0.7/0.10, 0.40/1.4/0.05,
    0.55/0.55/0.09, 0.68/0.9/0.07, 0.82/0.6/0.10, 0.93/0.75/0.08
  } {
    \fill[white, opacity=\o]
      ([xshift=1.4cm, yshift=-\y\paperheight]current page.north west)
      circle (\r cm);
  }

  \draw[white, line width=0.3pt, opacity=0.15]
    ([xshift=0.5cm]current page.north west) --
    ([xshift=0.5cm]current page.south west);
  \draw[white, line width=0.3pt, opacity=0.15]
    ([xshift=2.1cm]current page.north west) --
    ([xshift=2.1cm]current page.south west);
\end{tikzpicture}

% 第二层：顶部/底部装饰金线带
\begin{tikzpicture}[remember picture,overlay]
  \draw[GoldMain!45,line width=0.6pt]
    ([xshift=3.0cm,yshift=-1.6cm]current page.north west) --
    ([xshift=-2.0cm,yshift=-1.6cm]current page.north east);
  \foreach \x in {3.5,7,10.5,14} {
    \fill[GoldMain!55]
      ([xshift=\x cm,yshift=-1.45cm]current page.north west)
      circle (1.5pt);
  }
  \draw[GoldMain!35,line width=0.5pt]
    ([xshift=3.0cm,yshift=2.2cm]current page.south west) --
    ([xshift=-1.5cm,yshift=2.2cm]current page.south east);
  \foreach \x in {4,8,12} {
    \fill[GoldMain!40]
      ([xshift=\x cm,yshift=2.35cm]current page.south west)
      circle (1.2pt);
  }
\end{tikzpicture}

% 第三层：右上角装饰几何图形组（类波动同心圆环）
\begin{tikzpicture}[remember picture,overlay]
  \draw[GoldMain!30,line width=0.7pt]
    ([xshift=-2.5cm,yshift=-2.5cm]current page.north east)
    circle (22pt);
  \draw[WaveLight!25,line width=0.5pt]
    ([xshift=-2.5cm,yshift=-2.5cm]current page.north east)
    circle (14pt);
  \fill[GoldMain!50]
    ([xshift=-2.5cm,yshift=-2.5cm]current page.north east)
    circle (2.5pt);
  \foreach \ang in {0,90,180,270} {
    \draw[GoldMain!25,line width=0.4pt]
      ([xshift=-2.5cm,yshift=-2.5cm]current page.north east)
      -- ++(\ang:28pt);
  }
  \foreach \ang in {45,135,225,315} {
    \draw[WaveLight!18,line width=0.3pt]
      ([xshift=-2.5cm,yshift=-2.5cm]current page.north east)
      -- ++(\ang:26pt);
  }
\end{tikzpicture}

% 第四层：左下角装饰 —— 波动曲线簇
\begin{tikzpicture}[remember picture,overlay]
  \draw[GoldMain!15,line width=0.5pt,smooth]
    ([xshift=3.5cm,yshift=3.0cm]current page.south west)
    to[out=30,in=210] ++(3.0,2.0);
  \draw[WaveLight!12,line width=0.5pt,smooth]
    ([xshift=3.3cm,yshift=3.0cm]current page.south west)
    to[out=35,in=205] ++(3.2,2.2);
  \draw[GoldMain!12,line width=0.5pt,smooth]
    ([xshift=3.7cm,yshift=2.8cm]current page.south west)
    to[out=25,in=215] ++(2.8,1.8);
  \draw[WaveMain!10,line width=0.5pt,smooth]
    ([xshift=3.5cm,yshift=4.0cm]current page.south west)
    to[out=-15,in=195] ++(2.5,-1.5);
  \draw[GoldMain!10,line width=0.4pt,smooth]
    ([xshift=3.7cm,yshift=4.2cm]current page.south west)
    to[out=-20,in=190] ++(2.3,-1.3);
  \draw[GoldMain!22,-{Stealth[length=3pt,width=3pt]}]
    ([xshift=5.0cm,yshift=4.2cm]current page.south west)
    -- ++(0.8,0.5);
  \draw[WaveLight!18,-{Stealth[length=3pt,width=3pt]}]
    ([xshift=5.2cm,yshift=3.5cm]current page.south west)
    -- ++(0.7,-0.5);
\end{tikzpicture}

% 第五层：公式背景区 —— 淡金矩形 + 数学物理经典公式
\begin{tikzpicture}[remember picture,overlay]
  \shade[top color=GoldPale!40, bottom color=GoldPale!22]
    ([xshift=3.3cm,yshift=-3cm]current page.north west) rectangle
    ([xshift=-1.5cm,yshift=2.5cm]current page.south east);
  \draw[GoldMain!28,line width=0.4pt]
    ([xshift=3.3cm,yshift=-3cm]current page.north west) rectangle
    ([xshift=-1.5cm,yshift=2.5cm]current page.south east);

  % 七个经典公式
  \node[anchor=west,WaveMain!65] at ([xshift=3.7cm,yshift=-3.5cm]current page.north west)
    {$\displaystyle \oint_{C}\frac{f(z)}{z-a}\,\dd z
      =2\pi\iu\,f(a)$};

  \node[anchor=west,WaveLight!55] at ([xshift=3.8cm,yshift=-4.5cm]current page.north west)
    {$\displaystyle \oint_{C}f(z)\,\dd z
      =2\pi\iu\sum_{k}\Res f(z_k)$};

  \node[anchor=west,GoldMain!70] at ([xshift=3.9cm,yshift=-5.5cm]current page.north west)
    {$\displaystyle \frac{\partial^{2}u}{\partial t^{2}}
      =c^{2}\,\nabla^{2}u$};

  \node[anchor=west,WaveMain!50] at ([xshift=3.7cm,yshift=-6.5cm]current page.north west)
    {$\displaystyle \nabla^{2}u=0\qquad$};

  \node[anchor=west,WaveLight!45] at ([xshift=3.8cm,yshift=-7.5cm]current page.north west)
    {$\displaystyle x^{2}y''+xy'+(x^{2}-\nu^{2})y=0$};

  \node[anchor=west,GoldMain!55] at ([xshift=3.9cm,yshift=-8.5cm]current page.north west)
    {$\displaystyle \mathcal{F}[f](\xi)=\int_{-\infty}^{\infty}f(x)\,e^{-\iu\xi x}\,\dd x
      $};

  \node[anchor=west,WaveMain!42] at ([xshift=3.7cm,yshift=-9.5cm]current page.north west)
    {$\displaystyle \Gamma(z)=\int_{0}^{\infty}t^{z-1}e^{-t}\,\dd t$};

  \foreach \y/\c in {
    -3.5/WaveMain!30, -4.5/WaveLight!28, -5.5/GoldMain!32,
    -6.5/WaveMain!25, -7.5/WaveLight!22, -8.5/GoldMain!28, -9.5/WaveMain!20
  } {
    \fill[\c]
      ([xshift=-2.0cm,yshift=\y cm]current page.north west)
      circle (1.8pt);
  }
\end{tikzpicture}

% 第六层：标题区
\vspace*{2cm}

\hspace*{3.3cm}
\begin{minipage}{0.7\textwidth}
  {\large\bfseries\color{WaveLight} Huginn Notes in Mathematical Physics\par}

  \vspace{0.3cm}
  \begin{tikzpicture}
    \draw[GoldMain!60,line width=1.2pt] (0,0) -- (2.0,0);
    \fill[GoldMain!70] (2.3,0) circle (2pt);
    \fill[GoldMain!50] (2.8,0) circle (1.2pt);
    \fill[GoldMain!30] (3.2,0) circle (0.8pt);
  \end{tikzpicture}

  \vspace{0.8cm}

  {\Huge\bfseries\color{WaveDark} 数学物理方法笔记\par}

  \vspace{0.6cm}
  {\Large\color{WaveMain!70} Notes of Huginn Wong\par}

  \vspace{0.4cm}
  \begin{tikzpicture}
    \draw[GoldMain!40,line width=0.7pt] (0,0) -- (3.5,0);
    \fill[GoldMain!45] (0,0) circle (1.5pt);
    \fill[GoldMain!45] (3.5,0) circle (1.5pt);
  \end{tikzpicture}

  \vspace{0.8cm}
  {\large\color{GoldMain!80} $\mathfrak{For}~~\mathfrak{The}~~\mathfrak{Worthy}$\par}

\end{minipage}

% 第七层：右下角 logo
\begin{tikzpicture}[remember picture, overlay]
  \begin{scope}[shift={([xshift=-5.5cm, yshift=1.5cm]current page.south east)}]
    \draw[WaveMain, line width=0.8pt] (0,0) rectangle (0.9,0.72);
    \draw[WaveMain, line width=0.8pt] (0,0.72) -- (0.45,1.08) -- (0.9,0.72);
    \draw[GoldMain, line width=0.5pt] (0.45,0.72) -- (0.45,1.08);
    \draw[WaveMain, line width=0.4pt] (0.23,0) -- (0.23,0.72);
    \draw[WaveMain, line width=0.4pt] (0.67,0) -- (0.67,0.72);
    \fill[GoldMain] (0.45,1.08) circle (1.5pt);
    \node[anchor=west, WaveMain, font=\footnotesize] at (1.1,0.32)
      {Produced by Huginn};
  \end{scope}
\end{tikzpicture}

% 第八层：底部封面图片
\ifdefempty{\@coverimage}{}{
	\begin{tikzpicture}[remember picture,overlay]
		\node[
		anchor=south west,
		inner sep=0pt,
		outer sep=0pt
		] at ([xshift=3.3cm,yshift=3.05cm]current page.south west)
		{
			\includegraphics[
			width=\dimexpr\paperwidth-4.8cm\relax
			]{\@coverimage}
		};
	\end{tikzpicture}
}

% 第九层：散落装饰点
\begin{tikzpicture}[remember picture,overlay]
  \fill[GoldMain!18] ([xshift=-3.0cm,yshift=-9.5cm]current page.north west) circle (1.0pt);
  \fill[WaveLight!18] ([xshift=-3.8cm,yshift=-9.8cm]current page.north west) circle (0.7pt);
  \fill[GoldMain!16] ([xshift=-2.5cm,yshift=-10.2cm]current page.north west) circle (0.8pt);
  \fill[WaveMain!12] ([xshift=-4.0cm,yshift=3.5cm]current page.south west) circle (1.2pt);
  \fill[GoldMain!13] ([xshift=-5.5cm,yshift=4.0cm]current page.south west) circle (0.9pt);
  \fill[GoldMain!20] ([xshift=14.5cm,yshift=-6.0cm]current page.north west) circle (1.0pt);
  \fill[WaveLight!16] ([xshift=14.0cm,yshift=-5.5cm]current page.north west) circle (0.6pt);
  \fill[GoldMain!18] ([xshift=15.0cm,yshift=-5.8cm]current page.north west) circle (0.8pt);
\end{tikzpicture}

\end{titlepage}
}
\makeatother
